\section{模型检验与灵敏度分析}

本章对四个模型的关键结论作独立复核，并对高敏感假设做扰动分析，说明结果的适用范围。

\subsection{独立复核}
问题一另以不导入主评分模块的脚本从压缩原始记录重算，$51230$ 条 A1 逐条分数与保存值最大差异为 $2.3\times10^{-16}$，A2/A3 的域均值与冲突计数均一致。与 CRITIC、$22$ 项等权、argmax 解码的七域排序 Kendall 相关分别为 $0.333$、$0.238$、$0.810$，提示权重选择是主要不确定来源；A2/A3 的同名域均值接近 A1，不能证明其余域的泛化性。问题二的经典律在四组独立数据上分层验证；广义律在留出集的均方根误差 $0.688$ 低于同参经典基线 $0.971$，但 $Q=1$ 的零退化误差是公式恒等式。问题三把候选配置重代入物理 FLOPs 公式核验，三项总算力均未超预算；多起点与预算回代不能证明全局最优。问题四的一步滚动回测中，趋势模型平均绝对误差 $1.89$ 高于末值延续的 $1.28$。

\subsection{灵敏度分析}
质量放大指数 $\theta$ 不是可由现有质量标度识别的估计量，故对 $Q_0$ 与 $\theta$ 分别作条件扰动。在指数型质量成本、$\Lctx=2048$ 和同一多起点求解规则下，重求三档预算的候选配置；结果见表~\ref{tab:q3-assumption-sensitivity}。

\begin{table}[H]
  \centering
  \caption{质量基线与指数扰动下的候选最优质量}
  \label{tab:q3-assumption-sensitivity}
  \small
  \begin{tabular}{lrrrrr}
    \toprule
    情景 & $Q_0$ & $\theta$ & $Q^\ast_{10^{19}}$ & $Q^\ast_{10^{22}}$ & $Q^\ast_{10^{24}}$ \\
    \midrule
    基准 & 0.0789 & 0.5 & 0.546 & 0.968 & 1.000 \\
    $Q_0$ 降 20\% & 0.0631 & 0.5 & 0.546 & 0.968 & 1.000 \\
    $Q_0$ 升 20\% & 0.0947 & 0.5 & 0.545 & 0.968 & 1.000 \\
    $\theta=0.2$ & 0.0789 & 0.2 & 0.417 & 0.823 & 1.000 \\
    $\theta=0.8$ & 0.0789 & 0.8 & 0.623 & 1.000 & 1.000 \\
    \bottomrule
  \end{tabular}
\end{table}


$Q_0$ 上下扰动 $20\%$ 时，低预算候选质量从基准 $0.506$ 变为 $0.528$ 或 $0.585$；中预算约 $0.965$--$0.967$，高预算均为 $1$。因此低预算结果对新基线并非不敏感。固定 $Q_0$ 而将 $\theta$ 从 $0.2$ 改为 $0.8$，低预算质量由 $0.488$ 升至 $0.585$，中预算由 $0.821$ 升至边界 $1$。不同 $\theta$ 对应不同损失函数，情景间损失值不能解释为模型性能提升。

冲突降权因子 $\delta\in\{0.3,0.5,0.7,1.0\}$ 的扫描中七域排序均不变。三型质量成本在低预算的候选 $Q$ 为 $0.506$、$0.488$、$0.488$；中预算差异更大，为 $0.967$、$1$、$1$。能力前沿的三档斜率系数决定情景间差异，每档内的自助区间只覆盖趋势估计和月度残差，不包含情景选择或结构性模型误差。

\subsection{适用范围}
经典律的直接拟合支撑范围约为参数量 $0.07$--$12$、数据量至 $300$（均以 $10^9$ 计）；问题三的高预算解已超出该范围，其损失是公式外推。质量指数与能力前沿斜率系数均为高敏感假设，替代率、转移点及远期能力水平不宜作确定性解读。前沿仅有 $10$ 个有效月度点，短期回测也未证明趋势模型优于简单基线。
